\chapter*{Abstract}

\section*{\tituloPortadaEngVal}

The accumulation of content on contemporary digital platforms has surpassed the capacity of
classical search and filtering mechanisms to accompany the user in their exploration. Direct
search requires knowing in advance what is being sought, category-based filtering operates over
rigid taxonomies, and recommendation systems strip the user of agency by hiding the criteria
on which they operate. This Final Degree Project proposes a digital collection exploration
engine that combines three complementary mechanisms: faceted filtering over metadata,
transversal navigation between related collections through semantic references, and set
operations over partial results.

The system is structured as a client-server architecture. The navigation core has been
developed in Java over a relational database, exposing its capabilities through a
\textit{RESTful API}. The web interface, developed in JavaScript, provides the user with
interactive access to filters, links between collections, and the union set, which preserves
results saved during the session. Functional validation has been carried out over dumps of two
public catalogs of different nature, TMDB and DBLP, and the usability of the interface has
been evaluated through a formal user testing protocol.


\section*{Keywords}

\noindent Faceted navigation, collection exploration, hypermedia, relational filtering, set
operations, navigation engine, software engineering, usability evaluation, TMDB, DBLP.
